\section{Response to the R30 report}
We thank the referee for distinguishing the validity of the earlier verification arguments from the missing end-to-end numerical construction. The revised article replaces the action-independent numerical core with a prospectively specified action-dependent maintenance economy. Upper witnesses, own-policy values, admissible sets, and intervention costs are now constructed from primitives and charged. A new global theorem completes optimization over the unchanged active-boundary constant-control interval. We preserve the original economic target and all historical results. We do not claim neural acceleration or a solved all-domain feedback problem when the computed evidence does not establish them.

The source report is \path{reviews/2026-09-24-econometrica-r30/referee_report.md}, commit \path{e9bc144fbb6843d6a3436825a28584efb64fb8f1}. Its reviewed manuscript is commit \path{3394815f4ccbf5917582c8f756b535adf87783cc}. R32 branches from \path{6ea1e2fd8ab1928ded4ea98dc74ba6e1e72cf9d0}, retaining the R31 prospective protocol and amendments. This response addresses the report's individual findings rather than treating a universal final pass rate as closure of every scientific question.

\subsection{R30-F1: supplied witnesses versus a complete construction}
The main article's constructive theorem starts from rewards, terminal settlement, and action-dependent affine shock maps. Backward expectation and maximization produce $F_t$; shape-preserving dyadic compression produces $U_t$ with an exactly checked global error. The lower witness is explicitly $U_t-e_t$. Installed values are computed by exact own-policy recursion, not supplied. The work proposition includes these constructions and downstream cost-value partitions. Tables report witness time, intermediate and stored pieces, bit lengths, complete certificate bytes, and full cold-pipeline timing. This closes the missing input-construction step for the displayed coupled class; it does not establish that every such pipeline beats the strongest classical solver.

\subsection{R30-F2: cancellation of continuation values}
The new maintenance primitives are fixed in the inherited R31 protocol. Maintaining and replacing always have lower current reward than deferring, and have different transition maps. Their usefulness is entirely a continuation effect. The computation constructs the action-dependent Bellman alternatives at every date. No primitive reward-sign identity supplies the optimal policy. The old analytically reducible family is preserved in the historical annex, not used as the new coupled experiment.

\subsection{R30-F3: dominating classical controls}
The comparison now includes uncompressed exact piecewise-affine dynamic programming, an accuracy-only greedy solver using the same constructed upper witnesses, and 33/129-node spline fitted-Q proposals. These are more informative controls than a large grid. Their favorable performance is retained. The new contribution demonstrated by the experiment is the coupled construction and minimum-cost repair within its certified class. A neural end-to-end speed advantage is not established. We do not redefine an accuracy-only solver as a minimum-intervention solver, nor omit its much lower operating-solve time.

\subsection{R30-F4: policy-dependent occupancy}
The primary cost recursion uses the transition maps of the deployed action. A predeclared, separately labeled stress incumbent tests future exposure to an expensive installed instruction. In all six stress configurations, the dynamic solution has strictly smaller exact discounted intervention cost than the pointwise comparator using the same admissible sets. Both cost functions are evaluated under their own deployed occupancy. The other 36 cases show equality and are reported. Thus the previously absent mechanism is now exercised, without claiming it is universal or neural-specific.

\subsection{R30-F5: trillion-state rhetoric and dimension}
The new article makes no trillion-state headline comparison. Its geometric table reports state dimension, horizon, exact-value pieces, intermediate envelopes, stored witness pieces, and cost-value pieces. The earlier $2^{40}$ observation remains unchanged in the historical archive as a representation fact. The current study does not demonstrate high-dimensional scaling. Its output-sensitive construction theorem explicitly retains the possibility of growing partition complexity.

\subsection{R30-F6: the role of neural proposals}
The same frozen neural architecture, fitting recipe, and three seeds are evaluated at every declared horizon and tolerance. Raw neural passes are 0/18; the spline controls pass 12/12 and admit separate zero-intervention certificates. These findings are visible in the abstract, main table, and complete cohort. Neural proposals remain an evaluated installed representation, not a claimed enabling ingredient. The proposal-margin proposition states a verifiable condition under which an approximation eliminates revisions; no sampled training loss is substituted for its premise. Neural-specific superiority remains unestablished by these results.

\subsection{R30-F7: conditional complexity and proposal quality}
The upper-witness construction now has a mesh-and-precision existence bound and a complete output-sensitive work account. The cost of constructing the upper witness, own-policy evaluation, compilation, action sets, cost recursion, serialization, and checking is included. The proposal condition $\rho_t+2\kappa_t\le\eta$ quantitatively connects a verified approximation margin to eligibility and zero revision cost. The empirical study reports actual representation geometry rather than treating the favorable univariate case as a dimension-free theorem. No unproved speedup over a classical approximation is inferred.

\subsection{R30-F8: the original all-domain target}
The target remains exactly the original current-state, all-restart .01 operating guarantee. Its retained whole-domain upper bound is still 7.181834580823298, and the 47-coordinate derivative/residual bridge is not newly instantiated. We have not replaced this target with the maintenance tolerance. We make a substantive advance in the unchanged original economy through the global active-boundary scalar theorem, while keeping that smaller control class explicit. The article's constructive maintenance chain and the original continuous objective are now distinguished, rather than described as two parts of a completed feedback solve.

\subsection{R30-F9: complete optimization at the active boundary}
The five-point derivative diagnostic is replaced by an interval-wide proof for the same restart, controls, and parameter interval. The two drift half intervals have rigorously positive derivative lower bounds .065304 and .378008 after the upper-boundary error is charged. The unique global optimum in that constant class is $\theta=1/5$. Every Gaussian and exponential constant used in the proof is checked with rational arithmetic. The old validated endpoint payoff interval is identified as inherited; no new quadrature run is invented. This completes the scalar optimization requirement, but not unrestricted feedback optimization.

\subsection{R30-F10: arbitrary implementation prices}
The revised economic tables report the complete piecewise-linear envelope of $J_i-\lambda C_i$ over named computed deployments for every $\lambda\ge0$. All crossing prices and tie sets are stored as exact fractions. No 60/60 favorable sign at a chosen price is used as a headline. Cost units are normalized interventions, and no monetary calibration is asserted. The earlier illustrative-price experiment remains in the historical annex.

\subsection{R30-F11: the theorem-level distinction}
The article explicitly credits Bellman comparison and constrained backward induction as classical. The stated addition is a constructive, bit-controlled witness pipeline together with its generated action class and a checkable endogenous-cost representation. The cost recursion alone is not claimed as a new Bellman principle. A zero-intervention corollary also gives a global minimum over the true regret-feasible set when the computed incumbent certificate passes, avoiding unnecessary revisions caused by the conservative local class. A lower verified complexity than all established alternatives is not proved; the strong classical numerical comparisons are retained.

\subsection{R30-F12: one principal theorem-to-computation chain}
The article now proceeds from maintenance primitives to constructive witnesses, all-restart admissibility, endogenous minimum-cost revision, and exact economic sensitivity. The active-boundary result is a separate completed optimization theorem in the original economy. Earlier experiments and technical derivations are preserved in a lossless annex rather than interleaved as independent claims of superiority. The dominant chain is complete for the stated one-dimensional coupled class; a frontier-changing high-dimensional application is not claimed from it.

\section{Technical comments}
\paragraph{R30-T1: input complexity.} The charged-work proposition and stage accounting include witness production, raw-value evaluation, compilation, encoding, and checking. The empirical timing table additionally charges the independent checker.
\paragraph{R30-T2: constructing $U,L,v^0$.} $U$ is constructed by exact Bellman operations and global compression; $L=U-e$ is explicit; $v^0$ is an exact policy evaluation. None is loaded from the optimal reference.
\paragraph{R30-T3: online/offline amortization.} The article states the correct break-even formula including construction, loading, and per-query costs, and does not claim a measured break-even threshold because that deployment benchmark was not executed. The old exact online sign guard remains preserved as a strong control in its own action-independent model.
\paragraph{R30-T4: cardinality versus dimension.} Main complexity and evidence are expressed through actual one-dimensional partitions, horizon, bits, and bytes. No state-cardinality observation is used as evidence of dimension independence.
\paragraph{R30-T5: endogenous occupancy.} Six stress configurations show strict dynamic cost savings; 36 equality cases remain. Exact cost Bellman inequalities are independently checked, including isolated eligibility points.
\paragraph{R30-T6: proposal quality.} The verified sup-error/margin proposition connects approximation error to eligibility and zero revision cost. The raw/global own-policy test is reported separately from local-class repair. Unverified training residuals carry no certification force.
\paragraph{R30-T7: hardware execution.} Certificates concern exact compiled policies. Binary64 fitting coefficients are rationalized at compilation; actual floating-point neural forward execution is not certified.
\paragraph{R30-T8: economic units.} The weight $1+x$ defines normalized state-dependent implementation units. The complete shadow-price envelope is a sensitivity analysis, not an estimated institutional cost schedule.
\paragraph{R30-T9: boundary optimization.} The global monotonicity proof covers all $\theta\in[-1/5,1/5]$, includes the distant-boundary likelihood-score error, and yields the unique scalar optimum. It does not assume central-state concavity.
\paragraph{R30-T10: historical compression without deletion.} The current article is reorganized around the constructive coupled chain. \path{HISTORY_R32.pdf} reproduces all five R30 submission documents page-for-page, including their existing history. Earlier source and evidence paths are unchanged.

\section{Verification and remaining distinctions}
All 42 registered class certificates pass the independent checker. Twelve further zero-intervention certificates pass. Eight malformed-certificate tests are rejected, isolated-point/preimage tests pass, and a two-period exhaustive scenario-tree control agrees at four rational restarts. The scalar proof's rational inequalities all pass. These results support the specific mathematical and computational claims above. They are not a referee acceptance prediction, a proof of neural superiority, a high-dimensional performance study, or a completion of the original all-domain feedback target.
